% =================================================
% Учительская версия рабочего листа. Урок 1
% =================================================

\worksheettitle
  {Цифры, числа и разряды}
  {Урок 1 из 3. Версия для учителя}

\begin{teacherbox}[Цели урока]
  К концу урока ученик различает число и цифру, определяет разряд и значение
  цифры, понимает роль нуля в позиционной записи, раскладывает число по
  разрядам и восстанавливает число по разрядным слагаемым.
\end{teacherbox}

\begin{checkpointbox}[Входная диагностика: ответы]
  \begin{enumerate}[label=\textbf{\arabic*.},itemsep=0.35em]
    \item \(3\) цифры.
    \item Цифра \(8\).
    \item \(307=300+7\).
    \item \(4\,052\).
  \end{enumerate}
\end{checkpointbox}

\section{Из чего записывают числа}

Десять цифр:
\[
  0,\ 1,\ 2,\ 3,\ 4,\ 5,\ 6,\ 7,\ 8,\ 9.
\]

\begin{fillbox}[Заполненная таблица]
  \begin{center}
    \renewcommand{\arraystretch}{1.45}
    \begin{tabularx}{0.92\textwidth}{@{}>{\centering\arraybackslash}p{0.2\textwidth}>{\centering\arraybackslash}p{0.28\textwidth}X@{}}
      \toprule
      \rowcolor{TableHeader}
      Запись числа & Сколько цифр? & Цифры по порядку \\
      \midrule
      \(8\) & \(1\) & \(8\) \\
      \(47\) & \(2\) & \(4,7\) \\
      \(505\) & \(3\) & \(5,0,5\) \\
      \bottomrule
    \end{tabularx}
  \end{center}
\end{fillbox}

\Needspace{10\baselineskip}
\begin{teacherbox}[Как читать диагностику]
  Ошибка в задании 2 указывает на неустойчивое знание порядка разрядов.
  Запись \(307=30+7\) показывает, что ученик видит цифры, но не учитывает
  их позиции. Ответ \(452\) в задании 4 означает потерю нуля в разряде сотен.
\end{teacherbox}

\begin{propertybox}[Определения]
  \textbf{Цифра} --- знак, используемый для записи числа.
  \textbf{Число} показывает количество или порядок предметов.
\end{propertybox}

Однозначные числа в списке: \(7\) и \(9\). Многозначные: \(12\), \(304\),
\(6\,000\).

\begin{teacherbox}[Методический акцент]
  Не стоит требовать ответа «\(8\) --- это цифра или число?» без контекста.
  Знак \(8\) является цифрой, а запись \(8\) одновременно обозначает
  однозначное число. Лучше спрашивать: «Сколько цифр использовано в записи?»
\end{teacherbox}

\section{Одна цифра --- разные значения}

\begin{center}
  \renewcommand{\arraystretch}{1.45}
  \begin{tabularx}{0.94\textwidth}{@{}>{\centering\arraybackslash}p{0.16\textwidth}>{\centering\arraybackslash}p{0.32\textwidth}X@{}}
    \toprule
    \rowcolor{TableHeader}
    Число & Разряд цифры \(6\) & Значение цифры \(6\) \\
    \midrule
    \(6\) & единицы & \(6\) \\
    \(60\) & десятки & \(60\) \\
    \(600\) & сотни & \(600\) \\
    \(6\,000\) & единицы тысяч & \(6\,000\) \\
    \bottomrule
  \end{tabularx}
\end{center}

\begin{conclusionbox}
  Значение цифры зависит от её места, то есть от разряда.
\end{conclusionbox}

\section{Разряды многозначного числа}

\placevalueworksheetfigure

Для числа \(72\,649\):
\[
  7\rightarrow70\,000,\qquad
  2\rightarrow2\,000,\qquad
  6\rightarrow600,\qquad
  4\rightarrow40,\qquad
  9\rightarrow9.
\]

Цифры разрядов: десятки тысяч --- \(7\), единицы тысяч --- \(2\),
сотни --- \(6\), десятки --- \(4\), единицы --- \(9\).

\begin{teacherbox}[Устный вопрос]
  Спросите: «Что изменится, если цифры \(2\) и \(6\) поменять местами?»
  Важно добиться ответа: изменится не только порядок цифр, но и их значения;
  получится число \(76\,249\).
\end{teacherbox}

\section{Ноль сохраняет место разряда}

\zeroplacefigure

В числе \(40\,508\) цифра тысяч равна \(0\), цифра десятков равна \(0\).
После удаления нулей получится \(458\), и это другое число.

\begin{conclusionbox}
  Ноль внутри записи показывает, что единиц соответствующего разряда нет,
  и сохраняет места остальных цифр.
\end{conclusionbox}

\begin{teacherbox}[Типичная реплика ученика]
  «Ноль ничего не значит, его можно убрать». Предложите сравнить
  \(40\,508\) и \(458\), а затем назвать значение цифры \(4\) в каждом
  числе: \(40\,000\) и \(400\).
\end{teacherbox}

\section{Разложение числа по разрядам}

\expandedformfigure

\[
  47\,305=40\,000+7\,000+300+5.
\]

\begin{fillbox}[Ответы]
  \begin{enumerate}
    \item \(63\,204=60\,000+3\,000+200+4\).
    \item \(8\,090=8\,000+90\).
    \item \(70\,000+800+6=70\,806\).
    \item \(5\,000+40+3=5\,043\).
    \item \(40\,000+2\,000+50+9=42\,059\).
  \end{enumerate}
\end{fillbox}

\begin{warningbox}[Исправление ошибки]
  В записи \(50\,204=50\,000+2\,000+4\) цифру \(2\) ошибочно приняли
  за число тысяч. Она стоит в разряде сотен:
  \[
    50\,204=50\,000+200+4.
  \]
\end{warningbox}

\begin{teacherbox}[Проверка понимания, а не шаблона]
  После двух примеров попросите ученика объяснить, почему в разложении числа
  \(8\,090\) нет отдельного слагаемого для сотен и единиц. Ожидаемый ответ:
  в этих разрядах стоят нули.
\end{teacherbox}

\section{Самостоятельная проверка: ответы}

\begin{resultbox}
  \begin{enumerate}
    \item \(5\) цифр.
    \item Цифра \(3\).
    \item \(6\,000\).
    \item \(90\,307=90\,000+300+7\).
    \item \(34\,207\).
    \item \(40\,702\).
  \end{enumerate}
\end{resultbox}

\begin{checkpointbox}[Выходной билет: ответы]
  \begin{enumerate}
    \item Тремя цифрами.
    \item \(30\,408\).
    \item \(0\).
  \end{enumerate}
\end{checkpointbox}

\begin{teacherbox}[Критерий готовности к уроку 2]
  Ученик готов к изучению классов, если без подсказки:
  \begin{itemize}
    \item правильно движется по разрядам справа налево;
    \item сохраняет нули при восстановлении числа;
    \item различает цифру и значение цифры;
    \item выполняет хотя бы пять из шести заданий самостоятельной проверки.
  \end{itemize}
\end{teacherbox}
